% Working abstract. Journal limit: 150--250 words, no undefined abbreviations or references.
\begin{abstract}
\noindent Rodents live in acoustic environments where few sounds signal danger, yet failing to respond to the right ones can be fatal. How receivers assign meaning to sound, and why that meaning persists or fades with repetition, therefore bears directly on the function and evolution of antipredator behavior. We review how ground squirrels, marmots, and other field-studied rodents, together with laboratory mice and rats, recognize, interpret, and adapt to biologically meaningful sounds, including conspecific alarm and distress calls, heterospecific alarm calls, and predator vocalizations, relative to acoustically matched artificial controls. We organize the evidence around seven questions spanning recognition, assessment and behavioral response, and updating over time. Signal detection theory, threat-sensitive predator avoidance, the economics of escape, and the predation risk allocation hypothesis frame each question by asking when responding should be favored and when responsiveness should decline. We distinguish acoustic salience from ecological significance, reflexive startle from sustained defensive behavior, and sensory habituation from learning that a signal is unreliable, and we treat neural measures such as striatal dopamine release as evidence about proximate assessment rather than direct readouts of perceived threat. We close with a receiver-centered framework linking stimulus, prior experience, assessment, response, and updating, and propose testable hypotheses, pairing field and laboratory systems, about how signal reliability and temporal predictability govern the persistence of responses to repeated sounds.
\end{abstract}

\noindent\textbf{Keywords:} alarm calls; antipredator behavior; habituation; signal reliability; risk allocation; rodents
